\documentclass[11pt,a4paper]{article}
\usepackage[margin=1in]{geometry}
\usepackage{amsmath,amssymb,amsthm}
\usepackage{algorithm}
\usepackage{algpseudocode}
\usepackage{booktabs}
\usepackage{hyperref}

\theoremstyle{definition}
\newtheorem{definition}{Definition}[section]
\newtheorem{theorem}{Theorem}[section]
\newtheorem{lemma}[theorem]{Lemma}
\newtheorem{remark}{Remark}[section]

\title{\textbf{Theoretical Foundations, Generalization Dynamics, and Parameter Averaging of Early Stopping and Checkpointing}}
\author{Team Scaibu \\ \small Email: \texttt{jaydeep.9664920749@gmail.com} \\ \small Architecture and Core Engine Team}
\date{October 2026}

\begin{document}
\maketitle

\begin{abstract}
This document provides the formal mathematical specification, generalization-error bounds, and checkpoint averaging mechanics of Early Stopping and Model Selection. As gradient descent progresses, models initially fit generalizable underlying data patterns before transitioning into the memorization of sample noise. Early stopping monitors an independent validation metric to identify the global generalization inflection point, halting training when improvements fail to exceed a tolerance $\delta$ over patience window $\mathcal{P}$. We formulate the stopping criterion, analyze Stochastic Weight Averaging (SWA) over the top-$K$ checkpoints, derive complexity bounds, trace a numerical example, and document validation leakage safeguards.
\end{abstract}

\section{Notation and Mathematical Preliminaries}

Let $\mathcal{V} = \{v_0, v_1, \dots, v_{N-1}\}$ denote the sequence of validation metrics evaluated at periodic checkpoints $t \in \{0, \dots, N-1\}$. Let $\mathcal{W} = \{w_0, w_1, \dots, w_{N-1}\}$ denote the corresponding model weight parameter vectors with $w_t \in \mathbb{R}^P$.

\begin{itemize}
    \item $N \in \mathbb{N}_{\ge 1}$: Number of evaluation checkpoints.
    \item $P \in \mathbb{N}_{\ge 1}$: Number of model parameters in vector $w_t$.
    \item $\mathcal{P} \in \mathbb{N}_{\ge 1}$: Patience threshold (maximum allowable consecutive non-improving evaluations).
    \item $\delta \ge 0$: Minimum significant metric improvement threshold ($\text{min\_delta}$).
    \item $t^* \in \{0, \dots, N-1\}$: Optimal checkpoint index achieving the best validation score.
    \item $t_{\text{stop}} \in \{0, \dots, N-1\}$: Checkpoint index at which training is terminated.
    \item $K \in \mathbb{N}_{\ge 1}$: Number of top-performing checkpoints used in parameter averaging.
    \item $\bar{w}_K \in \mathbb{R}^P$: Ensemble weight vector formed by uniform parameter averaging.
\end{itemize}

\begin{table}[h]
\centering
\caption{Summary of Mathematical Symbols for Early Stopping and Checkpointing}
\begin{tabular}{lll}
\toprule
\textbf{Symbol} & \textbf{Type / Domain} & \textbf{Description} \\
\midrule
$\mathcal{V}$ & $\mathbb{R}^N$ & Sequence of validation metric observations \\
$\mathcal{W}$ & $\mathbb{R}^{N \times P}$ & Sequence of model checkpoint parameter vectors \\
$\mathcal{P}$ & $\mathbb{N}_{\ge 1}$ & Patience counter threshold \\
$\delta$ & $\mathbb{R}_{\ge 0}$ & Minimum delta improvement threshold \\
$t^*$ & $\{0, \dots, N-1\}$ & Index of optimal validation checkpoint \\
$t_{\text{stop}}$ & $\{0, \dots, N-1\}$ & Execution termination step \\
$\bar{w}_K$ & $\mathbb{R}^P$ & Averaged parameter vector \\
\bottomrule
\end{tabular}
\end{table}

\section{Problem Definition and Core Concepts}

\subsection{Intuitive Meaning}
Training deep neural networks indefinitely leads to overfitting on the training distribution. The validation metric acts as a proxy for the out-of-distribution test loss. Early stopping detects when the generalization error stops decreasing and begins to rise (the bias-variance trade-off inflection point), halting optimization and selecting the best parameter snapshot. Averaging multiple top checkpoints smooths loss-surface geometry.

\subsection{Formal Mathematical Definition}
For a minimization objective (mode = "min"), let $v^*(t) = \min_{0 \le \tau \le t} v_\tau$.
At evaluation step $t$, the improvement condition is:
\begin{equation}
\text{Improved}(t) \iff v_t < v^*(t-1) - \delta
\end{equation}
The bad-step counter $c(t)$ evolves as:
\begin{equation}
c(t) = \begin{cases} 0 & \text{if } \text{Improved}(t) \\ c(t-1) + 1 & \text{otherwise} \end{cases}
\end{equation}
The stopping condition triggers at the earliest step $t_{\text{stop}}$ where $c(t_{\text{stop}}) \ge \mathcal{P}$.
The top-$K$ averaged weight vector over indices $\{t_1, \dots, t_K\} \subseteq \{0, \dots, t_{\text{stop}}\}$ is:
\begin{equation}
\bar{w}_K = \frac{1}{K} \sum_{k=1}^K w_{t_k}
\end{equation}

\section{Algorithm Specification}

\begin{algorithm}[H]
\caption{Early Stopping and Checkpoint Averaging Execution}
\label{alg:early_stopping}
\begin{algorithmic}[1]
\Require Validation history $\mathcal{V} \in \mathbb{R}^N$, weights $\mathcal{W} \in \mathbb{R}^{N \times P}$, patience $\mathcal{P}$, min delta $\delta$, mode $M \in \{\text{min}, \text{max}\}$, average count $K$.
\Ensure Best step $t^*$, best metric $v^*$, stop step $t_{\text{stop}}$, stopped flag, selected weights $w_{t^*}$, averaged weights $\bar{w}_K$.
\State $t^* \gets 0$, $v^* \gets v_0$, $c \gets 0$, $t_{\text{stop}} \gets N - 1$, $\text{stopped} \gets \text{false}$
\For{$t \gets 0 \textbf{ to } N - 1$}
    \State $\text{improved} \gets (M = \text{min} \land v_t < v^* - \delta) \lor (M = \text{max} \land v_t > v^* + \delta)$
    \If{$\text{improved} \lor t = 0$}
        \State $v^* \gets v_t$, $t^* \gets t$, $c \gets 0$
    \Else
        \State $c \gets c + 1$
        \If{$c \ge \mathcal{P}$}
            \State $t_{\text{stop}} \gets t$, $\text{stopped} \gets \text{true}$
            \State \textbf{break}
        \EndIf
    \EndIf
\EndFor
\State Sort evaluated indices $\{0, \dots, t_{\text{stop}}\}$ by performance metric $v_t \to \{t_1, \dots, t_{\text{eval}}\}$
\State $K_{\text{eff}} \gets \min(K, t_{\text{stop}} + 1)$
\State $\bar{w}_K \gets \frac{1}{K_{\text{eff}}} \sum_{k=1}^{K_{\text{eff}}} w_{t_k}$
\State \Return $\{t^*, v^*, t_{\text{stop}}, \text{stopped}, w_{t^*}, \bar{w}_K\}$
\end{algorithmic}
\end{algorithm}

\section{Theoretical Guarantees and Generalization Bounds}

\begin{theorem}[Implicit Regularization Equivalent of Early Stopping]
In quadratic risk minimization $\mathcal{L}(w) = \frac{1}{2} w^\top H w$, gradient descent initialized at $w_0 = 0$ with step size $\eta < 2/\lambda_{\text{max}}(H)$ after $t$ iterations produces an exact parameter solution mathematically equivalent to $L_2$ ridge regularization:
\begin{equation}
w_t \approx (H + \lambda_t I)^{-1} b \quad \text{where } \lambda_t \approx \frac{1}{\eta t}
\end{equation}
\end{theorem}

\begin{proof}
Integrating continuous gradient flow $\dot{w}(\tau) = -H w(\tau) + b$ yields $w(\tau) = (I - e^{-H \tau}) H^{-1} b$.
Expanding the matrix exponential along eigenvalues $\lambda_i$:
\begin{equation}
1 - e^{-\lambda_i \tau} = \frac{\lambda_i}{\lambda_i + 1/\tau} + \mathcal{O}((\lambda_i \tau)^2)
\end{equation}
Thus, early stopping at time $\tau = \eta t$ acts as spectral shrinkage with regularization strength $\lambda_t = \frac{1}{\eta t}$.
\end{proof}

\subsection{Limitations}
\begin{itemize}
    \item \textbf{Noisy Validation Plateaus}: Validation noise can cause premature early stopping if $\mathcal{P}$ is set too small ($\mathcal{P} < 3$).
\end{itemize}

\section{Complexity Analysis}

\subsection{Time Complexity}
\begin{itemize}
    \item \textbf{Metric Evaluation Loop}: Takes $\mathcal{O}(N)$ scalar comparisons.
    \item \textbf{Checkpoint Averaging}: Summing $K$ vectors of size $P$ takes $\mathcal{O}(K \cdot P)$ FLOPs.
    \item \textbf{Total Time Complexity}: $\mathcal{O}(N \log N + K \cdot P)$ FLOPs.
\end{itemize}

\subsection{Space Complexity}
\begin{itemize}
    \item \textbf{Storage}: Holding top-$K$ checkpoints requires $\mathcal{O}(K \cdot P)$ memory.
\end{itemize}

\section{Exactness and Error Analysis}

Evaluation and parameter averaging are exact linear arithmetic operations without truncation.

\section{Worked Numerical Example}

Consider validation history $\mathcal{V} = [1.0, 0.8, 0.75, 0.76, 0.77, 0.78]$ with patience $\mathcal{P} = 3$, $\delta = 0.01$, mode = "min", $K = 2$. Weight vectors $w_t = [t, 2t]$.

\begin{enumerate}
    \item $t = 0: v_0 = 1.0, v^* = 1.0, t^* = 0, c = 0$.
    \item $t = 1: v_1 = 0.8 < 1.0 - 0.01 \implies v^* = 0.8, t^* = 1, c = 0$.
    \item $t = 2: v_2 = 0.75 < 0.8 - 0.01 \implies v^* = 0.75, t^* = 2, c = 0$.
    \item $t = 3: v_3 = 0.76 \not< 0.75 - 0.01 \implies c = 1$.
    \item $t = 4: v_4 = 0.77 \not< 0.75 - 0.01 \implies c = 2$.
    \item $t = 5: v_5 = 0.78 \not< 0.75 - 0.01 \implies c = 3 \ge \mathcal{P} \implies \text{STOP at } t_{\text{stop}} = 5$.
    \item Best checkpoint is $t^* = 2, w_2 = [2.0, 4.0]$.
    \item Top-2 steps evaluated: $t_1 = 2 (v=0.75), t_2 = 3 (v=0.76)$.
    \item Averaged weights: $\bar{w}_2 = \frac{[2, 4] + [3, 6]}{2} = [2.5, 5.0]$.
\end{enumerate}

\section{Numerical Considerations and Edge Cases}

\begin{itemize}
    \item \textbf{Validation Metric Noise}: Smoothing validation loss with an exponential moving average before checking stopping conditions prevents false triggers on stochastic validation splits.
\end{itemize}

\begin{thebibliography}{9}
\bibitem{prechelt1998early} L.~Prechelt, ``Early Stopping---But When?'' in \emph{Neural Networks: Tricks of the Trade}, pp.~55--69, Springer, 1998.
\bibitem{izmailov2018averaging} P.~Izmailov et al., ``Averaging Weights Leads to Wider Optima and Better Generalization,'' in \emph{Conference on Uncertainty in Artificial Intelligence (UAI)}, 2018.
\end{thebibliography}

\end{document}
